% -----------------------------------------------------------
\newpage
\subsection{Reglas de Negocio del Modelo}
% -----------------------------------------------------------

Las siguientes reglas definen las restricciones operativas y de consistencia que debe
respetar el modelo de datos. Su objetivo es preservar la integridad transaccional,
la trazabilidad histórica y la correcta separación entre operación en tiempo real,
liquidación interna y procesamiento analítico.

\begin{enumerate}
  \item \textbf{Herencia de configuración fiscal.}
  Si \texttt{Restaurant.taxBps} es \texttt{NULL}, se utiliza \texttt{Chain.defaultTaxBps}
  como tasa base. La resolución de herencia se realiza en la capa de aplicación antes
  de calcular totales o generar snapshots de orden.

  \item \textbf{Herencia de propinas sugeridas.}
  Si \texttt{Restaurant.suggestedTipsBps} es \texttt{NULL}, se utilizan las opciones
  definidas en \texttt{Chain.defaultTipsBps}. Estas propinas son únicamente sugerencias
  operativas y no implican procesamiento bancario por parte del sistema.

  \item \textbf{Eliminación lógica de catálogos.}
  Las entidades de catálogo no se eliminan físicamente cuando existe relación con
  registros históricos. En su lugar, se utiliza \texttt{archivedAt} o banderas equivalentes
  para preservar la integridad de órdenes, modificadores y métricas analíticas.

  \item \textbf{Snapshots históricos de órdenes.}
  Al crear un \texttt{OrderItem}, se copian los datos vigentes del producto, variante,
  estación, precio unitario, tasa de impuesto y nombre visible. Estos snapshots no deben
  modificarse aunque el catálogo cambie posteriormente.

  \item \textbf{Snapshots de modificadores.}
  Al seleccionar un modificador, \texttt{OrderItemModifier} almacena el nombre y ajuste
  de precio vigentes. Esto permite reconstruir el total histórico de una orden aun si el
  modificador original se archiva o cambia de precio.

  \item \textbf{Estados físicos de mesa.}
  \texttt{DiningTable.status} representa disponibilidad física-operativa y solo puede
  tomar los valores \texttt{LIBRE}, \texttt{OCUPADA} o \texttt{NO\_DISPONIBLE}. Una mesa
  física en estado \texttt{LIBRE} puede asociarse a una nueva sesión; una mesa
  \texttt{OCUPADA} no puede asociarse a otra sesión activa.

  \item \textbf{Estados de sesión virtual.}
  \texttt{DiningSession.status} representa el ciclo lógico de consumo. La transición
  normal es: \texttt{PENDIENTE\_ACTIVACION} $\rightarrow$ \texttt{ACTIVA}
  $\rightarrow$ \texttt{EN\_CONSUMO} $\rightarrow$ \texttt{POR\_LIQUIDAR}
  $\rightarrow$ \texttt{LIQUIDADA} $\rightarrow$ \texttt{CERRADA}. El estado
  \texttt{CANCELADA} se utiliza para anulaciones autorizadas.

  \item \textbf{Activación de sesión.}
  Una sesión inicia en \texttt{PENDIENTE\_ACTIVACION}. Pasa a \texttt{ACTIVA} cuando
  el anfitrión accede mediante QR firmado y define la contraseña de acceso, almacenada
  como \texttt{accessCodeHash}.

  \item \textbf{Inicio de consumo.}
  Una sesión pasa de \texttt{ACTIVA} a \texttt{EN\_CONSUMO} cuando se confirma la
  primera orden asociada a la sesión. La marca temporal correspondiente se registra en
  \texttt{firstOrderAt}.

  \item \textbf{Cierre operativo de mesa.}
  Al cerrar una \texttt{DiningSession}, las mesas físicas asociadas mediante
  \texttt{DiningSessionTable} vuelven a estado \texttt{LIBRE}, salvo que hayan sido
  marcadas administrativamente como \texttt{NO\_DISPONIBLE}.

  \item \textbf{Unión de mesas físicas.}
  Una sesión puede asociarse a una o varias mesas físicas mediante
  \texttt{DiningSessionTable}. Cada asociación registra \texttt{joinedAt} y, cuando la
  mesa deja de formar parte de la sesión, \texttt{leftAt}.

  \item \textbf{QR firmado y versionado.}
  El enlace QR se genera a partir de \texttt{DiningTable.publicCode} y
  \texttt{DiningTable.qrVersion}, firmados mediante HMAC-SHA256 del lado servidor.
  Incrementar \texttt{qrVersion} invalida enlaces anteriores sin modificar el identificador
  interno de la mesa.

  \item \textbf{Unicidad de identidad.}
  \texttt{AppUser.email} debe ser único en toda la plataforma. El alcance operativo del
  usuario se determina mediante \texttt{UserRole}, no mediante duplicación de cuentas.

  \item \textbf{Asignación contextual de roles.}
  Cada \texttt{UserRole} debe tener coherencia entre \texttt{contextType} y sus llaves
  contextuales. Si el contexto es \texttt{CHAIN}, debe existir \texttt{chainId}; si es
  \texttt{ZONE}, debe existir \texttt{zoneId}; si es \texttt{RESTAURANT}, debe existir
  \texttt{restaurantId}. Los permisos efectivos se calculan con base en el rol y el
  contexto asignado.

  \item \textbf{Órdenes y estados de cocina.}
  Los ítems de orden siguen el flujo \texttt{PENDIENTE} $\rightarrow$
  \texttt{EN\_PREPARACION} $\rightarrow$ \texttt{LISTA} $\rightarrow$
  \texttt{ENTREGADA}. Toda transición relevante debe registrarse en
  \texttt{OrderItemStatusEvent} para permitir trazabilidad y cálculo de tiempos.

  \item \textbf{Restricción durante liquidación.}
  Cuando una sesión se encuentra en \texttt{POR\_LIQUIDAR}, no se permiten nuevas órdenes
  ordinarias, salvo reapertura o autorización explícita por un usuario con permisos
  suficientes.

  \item \textbf{Liquidación interna sin pagos bancarios.}
  \texttt{Settlement} representa el estado interno de cuenta de una sesión. No almacena,
  transmite ni procesa datos bancarios. Las aportaciones se registran mediante
  \texttt{SettlementContribution} como eventos operativos internos.

  \item \textbf{Consistencia del saldo.}
  En todo momento debe cumplirse:
  \[
    \texttt{remainingCents} =
    \texttt{totalCents} - \texttt{amountSettledCents}
  \]
  considerando ajustes autorizados, impuestos y propina registrada internamente.

  \item \textbf{Aportaciones concurrentes.}
  El registro de aportaciones internas debe realizarse dentro de una transacción ACID.
  Para evitar sobreliquidación, el saldo pendiente se valida dentro de la transacción,
  utilizando bloqueo a nivel de fila sobre la liquidación correspondiente.

  \item \textbf{Asignación de aportaciones.}
  Una \texttt{SettlementContribution} puede cubrir una o varias asignaciones mediante
  \texttt{SettlementAllocation}. La suma de asignaciones asociadas a una contribución no
  debe exceder \texttt{SettlementContribution.amountCents}.

  \item \textbf{Cobertura de ítems.}
  La suma de \texttt{SettlementAllocation.amountCents} asociada a un
  \texttt{OrderItem} no debe exceder el total pendiente de dicho ítem.

  \item \textbf{Liquidación completa.}
  Una sesión solo puede pasar a \texttt{LIQUIDADA} cuando
  \texttt{Settlement.remainingCents = 0}. La fecha de liquidación se registra en
  \texttt{Settlement.settledAt} y \texttt{DiningSession.settledAt}.

  \item \textbf{Ajustes autorizados.}
  Los descuentos y cortesías se registran mediante \texttt{AppliedAdjustment}. Un ajuste
  puede aplicarse a un \texttt{OrderItem} o a un \texttt{Settlement}; cuando aplique una
  restricción XOR, no debe afectar ambos simultáneamente.

  \item \textbf{Anulación sin eliminación.}
  Las aportaciones y ajustes no se eliminan físicamente cuando requieren corrección. Se
  marcan como \texttt{VOIDED} y se registra el evento correspondiente en
  \texttt{AuditLog}.

  \item \textbf{Valores monetarios enteros.}
  Todos los montos se almacenan y calculan en centavos mediante enteros. No se utilizan
  tipos de punto flotante para dinero.

  \item \textbf{Auditoría de acciones críticas.}
  Toda acción crítica debe registrarse en \texttt{AuditLog}, incluyendo inicio de sesión,
  apertura o cierre de sesión de mesa, aplicación de ajustes, registro o anulación de
  aportaciones, cambios de estado relevantes y ejecución de jobs analíticos.

  \item \textbf{Separación OLTP y analítica.}
  Las consultas del dashboard histórico no deben recalcular métricas pesadas directamente
  sobre las tablas transaccionales. Deben consultar tablas Gold previamente generadas por
  jobs analíticos.

  \item \textbf{Control de jobs analíticos.}
  Cada ejecución del pipeline debe registrarse en \texttt{AnalyticsJobRun}, indicando
  tipo de job, estado, ventana de datos procesada, tiempos de inicio y fin, errores y
  metadatos técnicos.

  \item \textbf{Linaje de tablas Gold.}
  Cada registro en tablas Gold debe conservar \texttt{jobRunId}, permitiendo rastrear la
  métrica hasta la ejecución específica del pipeline que la generó.

  \item \textbf{Reprocesamiento incremental.}
  Los jobs analíticos deben utilizar \texttt{dataWindowStart} y
  \texttt{dataWindowEnd} para procesar ventanas específicas de datos. Esto permite
  actualizar métricas recientes o reprocesar rangos históricos sin recalcular todo el
  modelo.
\end{enumerate}

\newpage